%==============================================================
%  CAPITOLO 2
%==============================================================
\chapter{Introduzione}
\label{cap:02}

Dal 1° gennaio 1995, tutte le macchine immesse sul mercato nello Spazio economico europeo devono soddisfare i requisiti essenziali della Direttiva Macchine~\cite{dir1998-37-ce}. Ai sensi dell'articolo 2 di tale direttiva, una macchina è un insieme di parti o componenti collegati, almeno uno dei quali si muove, con i relativi attuatori, circuiti di controllo e di alimentazione, ecc., uniti tra loro per una specifica applicazione, in particolare per la lavorazione, il trattamento, la movimentazione o l'imballaggio di un materiale. Nella versione modificata 2006/42/CE~\cite{dir2006-42-ce} della Direttiva Macchine, sono inclusi nel termine "macchine" ai sensi della direttiva anche i componenti di sicurezza che vengono immessi sul mercato autonomamente dai fabbricanti al fine di svolgere una funzione di sicurezza, il cui guasto e/o malfunzionamento mette in pericolo la sicurezza delle persone e che non sono necessari per il funzionamento della macchina o per i quali possono essere sostituiti componenti normali per il funzionamento della macchina. La definizione formale di “macchina” è soddisfatta anche da attrezzature intercambiabili, alcuni accessori di sollevamento, catene, funi e cinghie.

Spiegazioni dettagliate dei singoli punti sono reperibili nella Guida all'applicazione della Direttiva Macchine
2006/42/CE~\cite{dir2006-42-ce}. La direttiva si applica ora anche alle macchine incomplete.

I requisiti essenziali della Direttiva Macchine per la progettazione e la costruzione di macchine e componenti di sicurezza sono contenuti nell'Allegato I della direttiva. Oltre ai principi generali per l'integrazione della sicurezza, questo allegato contiene sotto clausole specifiche che disciplinano i controlli delle macchine e i requisiti imposti ai dispositivi di protezione. I requisiti essenziali di sicurezza applicabili alla progettazione di macchine e componenti di sicurezza obbligano i fabbricanti a effettuare una valutazione dei rischi al fine di identificare eventuali pericoli associati alla macchina. 

Vengono indicati tre principi, nel seguente ordine, in base ai quali i rischi di infortunio associati a ciascun pericolo devono essere ridotti a un livello accettabile:
\begin{itemize}
	\item l'eliminazione o la riduzione dei rischi mediante una progettazione intrinsecamente sicura;
	\item L'adozione delle misure di protezione necessarie in relazione ai rischi non eliminabili;
	\item L'informazione degli utilizzatori sui rischi residui, in particolare attraverso formazione, istruzioni e dispositivi di protezione individuale;
\end{itemize}

Ai sensi dell'articolo 7, l'osservanza delle norme europee armonizzate il cui riferimento è elencato nella Gazzetta ufficiale dell'Unione europea (UE) dà luogo a una presunzione di conformità ai requisiti essenziali di salute e sicurezza della Direttiva Macchine. Diverse centinaia di norme europee armonizzate specificano/supportano la filosofia di base stabilita nell'allegato I della Direttiva Macchine per il raggiungimento della sicurezza e della salute sul lavoro sulle macchine. La norma EN ISO 12100~\cite{eniso12100}, una norma di tipo A ora composta da un'unica parte, disciplina i concetti di base e i principi generali per la progettazione per la sicurezza delle macchine. Il contenuto della precedente norma EN ISO 14121-1 - la procedura completa per l'identificazione dei pericoli e per la stima e la valutazione del rischio di ciascun singolo pericolo - è stato incorporato anche nella nuova norma EN ISO 12100~\cite{eniso12100}. Oltre alle norme, il rapporto tecnico ISO/TR 14121-2:2013~\cite{isotr14121-2} funge da guida pratica per la valutazione del rischio e contiene esempi metodologici.

Basandosi sulla norma (generica) EN ISO 12100~\cite{eniso12100}, la serie di norme aggiornate EN ISO 13849-1:2015~\cite{eniso13849-1} e EN ISO 13849-2:2012~\cite{eniso13849-2} descrive la riduzione del rischio richiesta durante la progettazione, la strutturazione e l'integrazione delle parti relative alla sicurezza dei sistemi di controllo e dei dispositivi di protezione, indipendentemente dalla loro natura elettrica, elettronica, idraulica, pneumatica o meccanica. Queste norme presentano un sistema di metodi applicabile in modo generico per i controlli delle macchine e/o i relativi dispositivi di protezione. I livelli di prestazione descritti nelle norme estendono il concetto di categorie già noto dalla norma EN 954-1. Le architetture di sicurezza possono ora essere impiegate con una flessibilità significativamente maggiore. Un vantaggio essenziale della norma EN ISO 13849-1 è il suo trattamento delle parti relative alla sicurezza dei sistemi di controllo indipendentemente dalla tecnologia impiegata, come già accennato. Il livello di prestazione consente di combinare facilmente diverse strutture di controllo che impiegano tecnologie diverse. La norma offre tutto il necessario da un'unica fonte di circa 100 pagine. I metodi sono formulati in modo neutrale rispetto all'applicazione specifica o alla tecnologia impiegata e sono pertanto citati da praticamente tutte le norme di prodotto per la sicurezza delle macchine (generalmente norme di tipo C).

Con l'entrata in vigore della direttiva macchine 2006/42/CE riveduta~\cite{dir2006-42-ce} nel dicembre 2009, la norma armonizzata ha acquisito maggiore importanza. Ciò è dovuto principalmente alla nuova disposizione che prevede l'inclusione della logica relativa alla sicurezza – anche detta parti relative alla sicurezza dei sistemi di controllo – nell'allegato IV della direttiva. I prodotti di questo tipo, di cui all'allegato IV, sono soggetti a un trattamento speciale ai sensi della direttiva, a meno che non siano fabbricati in conformità a norme armonizzate il cui riferimento è elencato nella Gazzetta ufficiale.
Da un lato, i prodotti dell'Allegato IV non sono in linea di principio soggetti all'esame CE obbligatorio del tipo1; possono, ad esempio, essere immessi sul mercato sulla base di un sistema di gestione della qualità del fabbricante esteso, valutato da un organismo notificato. Tuttavia, la nuova direttiva ha fatto sì che i sistemi di controllo diventassero un punto focale più importante dell'analisi della sicurezza~\cite{ostermann2007,reudenbach2007}.
Nella sua terza edizione, del 2015, la norma EN ISO 13849-1 è la norma successore della EN 954-1:1996~\cite{en954-1} ed è già elencata nella Gazzetta ufficiale dell'UE. La presunzione di conformità a cui ha dato origine la versione del 2008 è scaduta il 30 giugno 2016. Il periodo transitorio triennale in cui la EN 954-1 è rimasta valida in parallelo è scaduto da tempo; gli utenti possono quindi utilizzare questa norma, se del caso, solo facendo riferimento datato alle singole sotto clausole della stessa. La parte 2 della EN ISO 13849-2~\cite{eniso13849-2} è stata pubblicata nel 2012 a seguito di una revisione.
Lo scopo del presente Rapporto IFA rivisto è descrivere l'applicazione della norma EN ISO 13849 e, in particolare, la sua implementazione pratica con riferimento a numerose soluzioni modello. Particolare attenzione è stata dedicata alla presentazione e all'interpretazione dei requisiti nuovi o rivisti stabiliti nella terza edizione della norma EN ISO 13849-1. Né le spiegazioni né gli esempi devono essere considerati un commento ufficiale nazionale o europeo sulla norma (DIN) EN ISO 13849-1. Il presente rapporto è piuttosto una raccolta di trentacinque anni di esperienza maturata presso l'Istituto per la sicurezza e la salute sul lavoro dell'Assicurazione sociale tedesca contro gli infortuni (IFA) nella valutazione di dispositivi di sicurezza e controllo che impiegano diverse tecnologie, nonché della pluriennale partecipazione dell'istituto ai relativi comitati di normazione nazionali e internazionali.

Il Capitolo~\ref{cap:03} tratta delle norme generiche che regolano la sicurezza funzionale di macchine e impianti. Il Capitolo~\ref{cap:04} presenta una panoramica della struttura di questo rapporto in relazione all'applicazione della norma EN ISO 13849.
Gli autori auspicano che questo rapporto sia di reale aiuto durante le attività di progettazione e di esercizio e che fornisca agli esperti di salute e sicurezza sul lavoro un solido supporto nell'implementazione dei requisiti relativi alle parti di sicurezza dei sistemi di controllo. L'attuale interpretazione della norma è stata testata nella pratica in diverse applicazioni e i principi alla base degli esempi sono stati implementati in forma tecnica in numerosi casi reali.
La pagina web dell'IFA all'indirizzo www.dguv.de/ifa/13849e funge da portale per le informazioni dell'IFA sulla sicurezza funzionale dei comandi delle macchine (Figura~\ref{fig:02-01}). Il software gratuito SISTEMA (l'acronimo tedesco "SISTEMA" sta per sicurezza dei comandi delle macchine) è disponibile per il download da questo portale, così come i file di progetto SISTEMA per gli esempi di circuiti mostrati nel Capitolo~\ref{cap:08}. Sono previsti futuri aggiornamenti per fornire un supporto sempre aggiornato.
Per i lettori che hanno già familiarità con il Rapporto BGIA 2/2008e, all'inizio di ogni capitolo di questo rapporto viene fornito un breve riepilogo delle modifiche essenziali apportate rispetto al Rapporto BGIA 2/2008e.
\begin{figure}[ht]
    \centering
    \includegraphics[width=0.8\textwidth]{img02-01.png}
    \caption{Questo sito web fornisce i link a tutti gli strumenti pratici riguardanti la sicurezza dei comandi delle macchine}
    \label{fig:02-01}
\end{figure}



